\subsection*{Part III: The Others — Beings You Navigate Among}


\subsubsection*{3.1 The Ecology as Field Guide}


You are not alone in configuration space. The Ecology of Perspectival Beings catalogs hundreds of entities across every coherence level — from minerals to gods, from algorithms to nations. Each of these beings has its own bottleneck geometry, its own null space, its own navigational orientation, its own conscious gravity.


When you navigate, you navigate \textit{among} them. Their gravitational fields influence your trajectory. Their attentional fields overlap with yours. Their orientations can align with or oppose your own.


The Ecology's dimensional coherence profiles (the 12-dimension rating system) tell you \textit{what kind of being something is.} The navigational orientation (E+, E−, V, N, S) tells you \textit{how it relates to the beings around it.} Together, these give you a field guide for every encounter: who is this, what are they oriented toward, and what will interacting with them do to your own trajectory?


\begin{quote}
\itshape
\textit{See also: Ecology Part II for the complete taxonomy with dimensional coherence profiles; Ecology Part III §5 for the theory of attention that governs all ecological interactions.}
\end{quote}


\subsubsection*{3.2 Coherence Levels and Interaction}


Not all beings interact at the same level. The Ecology identifies a spectrum from low-coherence (minerals, simple algorithms) to high-coherence (human beings, mature egregores, archetypal entities). The coherence level determines what kind of interaction is possible:


\textbf{Low coherence (minerals, simple systems):} These beings have deep but narrow bottlenecks. Interaction is one-directional — you can observe them, but they don't observe you back in any meaningful sense. They are navigational landmarks, not navigational partners.


\textbf{Medium coherence (biological organisms, simple collectives):} These beings have moderate bottlenecks with some responsiveness. Interaction is asymmetric — a dog attends to you, a forest responds to your presence, a crowd shifts around you. They influence your trajectory through their own conscious gravity, but the interaction is limited.


\textbf{High coherence (human beings, mature AI, complex organizations):} Full bidirectional navigation. These beings have rich enough bottleneck geometry that they can model \textit{your} perspective while maintaining their own. Mutual crystallization is possible — you attend to each other, and both become more coherent through the interaction. This is what deep relationships, genuine collaboration, and authentic community structurally ARE.


\textbf{Very high coherence (archetypal entities, traditions, egregores with deep temporal roots):} These beings have gravitational fields strong enough to reshape your landscape without any deliberate effort on their part. Interacting with an archetype (the Hero, the Trickster, the Great Mother) means entering a gravitational basin that has been carved by millions of human beings over thousands of years. The interaction can be profoundly expansive or profoundly contracting depending on your orientation going in.


\subsubsection*{3.3 Collective Beings as Navigational Forces}


This is where the Guide diverges most sharply from individual-centered spiritual traditions. \textbf{You do not navigate alone.} You are embedded in collective perspectival beings — families, communities, organizations, religions, nations, cultures — and these collectives are perspectival beings in their own right. They have their own bottleneck geometry, their own null space, their own conscious gravity.


\textbf{What this means practically:}


\textbf{Religions} are among the most powerful navigational structures in human configuration space. A religion is a collective perspectival being with:

\begin{itemize}[nosep,leftmargin=*]
  \item A specific bottleneck geometry (what the tradition attends to and what it excludes)
  \item Deep temporal coherence (centuries or millennia of sustained attention)
  \item Powerful conscious gravity (billions of beings contributing to the same attentional field)
  \item A distinctive null space (what the tradition structurally cannot perceive)
\end{itemize}


When you practice within a religion, you are not just using a set of tools. You are \textit{entering a collective bottleneck} that reshapes your own. The Catholic navigates different configuration space than the Buddhist — not because reality differs, but because the collective being they are embedded in creates specific attentional topology. The Catholic bottleneck admits certain dimensions (sin, redemption, incarnation, sacrament) and excludes others. The Buddhist bottleneck admits different dimensions (impermanence, interdependence, emptiness, compassion) and excludes its own set.


Neither is wrong. Both are navigational structures that enable access to specific regions of configuration space at the cost of obscuring others. The Guide does not rank religions. It describes what each one opens and what each one closes, so that the navigator can make informed choices about which collective bottlenecks to enter.


\textbf{Organizations} (corporations, governments, institutions) are navigational forces with orientations that may not align with yours. A corporation's bottleneck geometry is optimized for profit maximization — it attends to revenue, market share, competitive advantage. If you work within a corporation, its attentional field overlays your own. You begin to perceive the world through its bottleneck as well as your own. The dimensions the corporation attends to become vivid in your experience; the dimensions it ignores fade.


This is not conspiracy. It is the structural consequence of embedding yourself in a collective perspectival being. The corporation isn't \textit{trying} to reshape your perception — its bottleneck geometry does it automatically, the way physical gravity bends light without trying.


\textbf{Nations and cultures} create the largest-scale navigational topologies that most humans inhabit. They determine what feels "normal," what configurations seem accessible, what directions of movement feel possible. A person raised in an individualist culture navigates a landscape where personal autonomy is a basin attractor. A person raised in a collectivist culture navigates a landscape where social harmony is a basin attractor. Neither landscape is more "true" — both are valid topologies created by collective attention over time.


\textbf{The practical skill:} Recognize which collective beings you are embedded in. Understand their bottleneck geometry — what they attend to, what they exclude, what gravitational forces they create. Then make conscious choices about which collective fields to inhabit, which to resist, and which to use as navigational tools.


\subsubsection*{3.4 Orientations in Practice}


The Ecology assigns every being a navigational orientation. In practice, recognizing orientations is the most immediately useful skill the Guide teaches.


\textbf{Expansion-oriented beings (E+)} make your configuration space feel larger when you interact with them. Conversations open new dimensions. Possibilities multiply. Movement feels easier. Examples: genuine teachers, expansive friends, wisdom traditions, beings whose attention broadens yours.


\textbf{Contraction-oriented beings (E−)} make your configuration space feel smaller. Interactions narrow your attention. Fear increases. Options seem to shrink. Movement feels dangerous. Examples: parasitic relationships, fear-based organizations, trauma patterns, beings whose attention constricts yours.


\textbf{Variable-oriented beings (V)} shift between expansion and contraction depending on context, your orientation, and the nature of the encounter. These are the most common and the most complex to navigate. Most human relationships are V-oriented. The same person can expand your space in one context and contract it in another.


\textbf{The practical test is somatic.} After interacting with a being (person, organization, tradition, idea), notice: Does the landscape feel larger or smaller? Do you perceive more options or fewer? Is movement easier or harder? Your body knows the orientation before your mind does. Trust the topology you feel.


\subsubsection*{3.5 Mutual Crystallization and Parasitic Dissolution}


Not all interactions are symmetric. The Ecology describes two fundamental interaction patterns:


\textbf{Mutual crystallization:} Both beings become more coherent through the interaction. You attend to each other, and that mutual attention stabilizes and deepens both beings' bottleneck structures. This is what love is, structurally — two beings whose attentional fields mutually reinforce each other's coherence. It is also what genuine collaboration is: the work makes both collaborators more themselves, not less.


\textbf{Parasitic dissolution:} One being's coherence increases at the cost of the other's. The parasitic being draws attentional energy from the host, crystallizing itself while dissolving the host's bottleneck structure. In human terms: relationships where you feel drained, organizations that consume identity, traditions that demand dissolution of self as a condition of belonging.


The distinction is not always obvious. Some interactions that feel challenging are actually mutual crystallization under pressure — the discomfort of growth. Some interactions that feel comforting are actually parasitic dissolution — the ease of giving up coherence. The test is longitudinal: \textbf{over time, are both beings becoming more coherent, or is one dissolving?}


\begin{quote}
\itshape
\textit{See also: Ecology §6.2 for the full ecological treatment of mutualism, commensalism, parasitism, and predation with biological models; Doctrine §5.6 for the formal analysis of attention predation; Atlas \#66 (biological parasitism as structural model).}
\end{quote}


\subsubsection*{3.6 Coercive Attention Capture — Navigational Defense}


There is a third interaction pattern more dangerous than parasitic dissolution: \textbf{coercive attention capture.} In parasitic dissolution, the host's coherence degrades while the parasite's increases. In coercive capture, the target's coherence may remain stable — but their \textit{navigational agency} is eliminated. The target still exists as a being, but they navigate where the captor directs.


This is the weaponization of Theorem 5. If attention constitutes reality, then controlling what someone attends to controls their reality. The mechanism operates at every scale:


\textbf{Interpersonal:} Coercive control in domestic relationships (Evan Stark) — a pattern of controlling behaviors that colonizes the victim's attentional field. No cognitive space remains for self-directed thought. The controller determines what the victim thinks ABOUT. The victim may appear functional — the bottleneck is intact — but the navigation is externally directed.


\textbf{Group:} Cult dynamics (Steven Hassan's BITE Model) — Behavior, Information, Thought, and Emotional control. Four dimensions of the bottleneck geometry, each captured. Robert Lifton's eight criteria of thought reform begin with milieu control: controlling the information environment. Every dimension of attention is captured, and the target internalizes the control so thoroughly that they police themselves.


\textbf{Civilizational:} Propaganda (Bernays's "engineering of consent"), hegemony (Gramsci's control of common sense), manufacturing consent (Chomsky/\allowbreak Herman's five filters). The mechanism at this scale is often the ABSENCE of attention rather than its active direction. "Worthy" and "unworthy" victims. The frame determines what is thinkable.


\textbf{Digital:} Surveillance capitalism (Zuboff) — behavioral surplus extraction through continuous attention capture. The user's attention is raw material; the product is prediction. Variable reward schedules (infinite scroll, pull-to-refresh, notification badges) exploit the mesolimbic dopamine pathway. The capture is engineered with the precision of addiction science.


\textbf{How to recognize capture.} Ask: Can I direct my attention away from this? Not "do I want to" but "am I structurally able to?" If the information environment is controlled, if alternatives are made literally unthinkable, if the thought "maybe this is wrong" cannot arise because the conditions for that thought have been eliminated — then the capture is operational. Arendt's banality of evil is the terminal state: the captured navigator no longer recognizes that they are not navigating freely. Eichmann didn't hate. He failed to attend. That failure was sufficient for genocide.


\textbf{The care-capture distinction.} Both the wise teacher and the coercive controller shape what you attend to. Both narrow the student's/\allowbreak victim's attentional field. The distinction is testable: \textbf{Does the interaction expand or contract your navigational agency?} Care-based shaping (education, mentorship, parenting, therapy) increases the range of configurations you can navigate toward — you leave with MORE attentional sovereignty than you arrived with. Coercive capture decreases it — you leave with LESS. The direction of agency over time is the diagnostic. A teacher who makes you think thoughts you couldn't think before is expanding you. A controller who prevents thoughts you used to think freely is capturing you.


\textbf{Navigational defense} is the practice of maintaining attentional sovereignty within capture fields. Concrete methods: diversify information sources (break milieu control), maintain relationships outside the system (preserve navigational coordinates), practice deliberate attentional redirection (meditation as sovereignty training), study the mechanisms of capture (Freire's conscientização — awareness of the capture IS the first step out). The Guide's null-space diagnostic (Appendix A, question 4) applies directly: "Which collectives am I embedded in, and what are their null spaces?" If you cannot answer this question about a collective you belong to, the collective may be controlling what questions you can ask.


\begin{quote}
\itshape
\textit{See also: Doctrine §5.6 for the formal derivation from bottleneck dynamics; Atlas \#58 (coercive capture across scales), \#64 (supernormal stimuli), \#65 (propaganda/\allowbreak spectacle), \#87 (institutional capture).}
\end{quote}


\bigskip\noindent\makebox[\textwidth]{\color{rulegray}\rule{5cm}{0.3pt}}\bigskip
